\documentclass[11pt,a4paper,oneside]{article}

% =====================================================
% Configuración central: portada, estilo y bibliografía
% =====================================================
\usepackage[utf8]{inputenc}
\usepackage[T1]{fontenc}
\usepackage[spanish,es-nodecimaldot]{babel}
\addto\captionsspanish{\renewcommand{\tablename}{Tabla}}
\usepackage{newtxtext,newtxmath}
\renewcommand{\ttdefault}{cmtt}
\usepackage[a4paper,left=3cm,right=3cm,top=2.5cm,bottom=2.5cm]{geometry}
\usepackage{setspace}
\onehalfspacing
\usepackage{graphicx,xcolor}
\usepackage{pdflscape}
\usepackage{capt-of}
\graphicspath{{docs/figuras/}{figuras/}}
\usepackage{booktabs,longtable,tabularx,array}
\usepackage{csquotes}
\usepackage{xurl}
\usepackage{hyperref}
\usepackage[backend=biber,style=ieee,citestyle=ieee,sorting=none]{biblatex}
\setcounter{biburlnumpenalty}{100}
\setcounter{biburlucpenalty}{100}
\setcounter{biburllcpenalty}{100}

% Permite compilar desde la raíz del repositorio o desde docs/.
\IfFileExists{docs/referencias/referencias.bib}{%
  \addbibresource{docs/referencias/referencias.bib}%
}{%
  \addbibresource{referencias/referencias.bib}%
}
\newcommand{\seccion}[1]{%
  \IfFileExists{docs/secciones/#1.tex}{%
    \input{docs/secciones/#1.tex}%
  }{%
    \input{secciones/#1.tex}%
  }%
}
\hypersetup{colorlinks=true,linkcolor=black,citecolor=black,urlcolor=blue!55!black}
\urlstyle{same}
\newcommand{\pendiente}[1]{{\color{red!65!black}\textbf{[Pendiente: #1]}}}

% Datos compartidos por la portada del informe y la carátula del SGA.
\IfFileExists{docs/datos-portada.tex}{% Datos compartidos por el informe y la carátula de entrega.
\newcommand{\nombreproyecto}{Sistema de Gestión de Refuerzo Académico (SGRA)}
\newcommand{\asignatura}{Aplicaciones Distribuidas}
\newcommand{\periodoacademico}{REGULAR 2026--2027 SPA}
\newcommand{\fechadeentrega}{21 de noviembre de 2026}
\newcommand{\repositoriopfc}{https://github.com/Yukisan-uwo/PFC-SGRA}
\newcommand{\arquitecto}{LOMBEIDA ESCALERAS BRYAN HUMBERTO}
\newcommand{\liderdesarrollo}{CORONEL GOMEZ KHRIZ ALEXANDER}
\newcommand{\responsablecalidad}{ZAMBRANO PARRAGA LUIS MARIO}
\newcommand{\responsabledocumentacion}{TOAQUIZA ZAMBRANO BELINDA BETZABETH}

\newcommand{\facultad}{Facultad de Ciencias de la Computación}
\newcommand{\carrera}{Ingeniería de Software}
\newcommand{\curso}{Séptimo semestre, paralelo A}
}{% Datos compartidos por el informe y la carátula de entrega.
\newcommand{\nombreproyecto}{Sistema de Gestión de Refuerzo Académico (SGRA)}
\newcommand{\asignatura}{Aplicaciones Distribuidas}
\newcommand{\periodoacademico}{REGULAR 2026--2027 SPA}
\newcommand{\fechadeentrega}{21 de noviembre de 2026}
\newcommand{\repositoriopfc}{https://github.com/Yukisan-uwo/PFC-SGRA}
\newcommand{\arquitecto}{LOMBEIDA ESCALERAS BRYAN HUMBERTO}
\newcommand{\liderdesarrollo}{CORONEL GOMEZ KHRIZ ALEXANDER}
\newcommand{\responsablecalidad}{ZAMBRANO PARRAGA LUIS MARIO}
\newcommand{\responsabledocumentacion}{TOAQUIZA ZAMBRANO BELINDA BETZABETH}

\newcommand{\facultad}{Facultad de Ciencias de la Computación}
\newcommand{\carrera}{Ingeniería de Software}
\newcommand{\curso}{Séptimo semestre, paralelo A}
}

\begin{document}

% =====================================================
% 1. PORTADA OFICIAL
% =====================================================
\IfFileExists{docs/portada.tex}{\begin{titlepage}
\begin{center}
  {\Large\bfseries Universidad Técnica Estatal de Quevedo}\par
  \vspace{0.4cm}
  {\large\bfseries \facultad}\par
  \vspace{0.8cm}
  {\large\bfseries \asignatura}\par
  \carrera\par
  \curso\par
  \vspace{0.8cm}
  {\Large\bfseries Proyecto Fin de Curso --- Entrega 1}\par
  \vspace{0.4cm}
  {\LARGE\bfseries \nombreproyecto}\par
  \vspace{0.4cm}
  {\large TA-PFC-E1: Propuesta y Análisis del Sistema Distribuido}\par
\end{center}
\vspace{0.8cm}
\noindent\textbf{Integrantes y roles:}\par\smallskip
\noindent\begin{tabularx}{\textwidth}{@{}p{0.34\textwidth}X@{}}
Arquitecto: & \arquitecto\\[0.15cm]
Líder de Desarrollo: & \liderdesarrollo\\[0.15cm]
Responsable de Calidad: & \responsablecalidad\\[0.15cm]
Responsable de Documentación: & \responsabledocumentacion\\
\end{tabularx}
\vfill
\begin{center}
\textbf{Período académico:} \periodoacademico\par
\medskip
\textbf{Fecha de entrega:} \fechadeentrega\par
\bigskip
\textbf{Repositorio del proyecto:}\par
\smallskip
\mbox{\url{\repositoriopfc}}
\end{center}
\end{titlepage}
}{\begin{titlepage}
\begin{center}
  {\Large\bfseries Universidad Técnica Estatal de Quevedo}\par
  \vspace{0.4cm}
  {\large\bfseries \facultad}\par
  \vspace{0.8cm}
  {\large\bfseries \asignatura}\par
  \carrera\par
  \curso\par
  \vspace{0.8cm}
  {\Large\bfseries Proyecto Fin de Curso --- Entrega 1}\par
  \vspace{0.4cm}
  {\LARGE\bfseries \nombreproyecto}\par
  \vspace{0.4cm}
  {\large TA-PFC-E1: Propuesta y Análisis del Sistema Distribuido}\par
\end{center}
\vspace{0.8cm}
\noindent\textbf{Integrantes y roles:}\par\smallskip
\noindent\begin{tabularx}{\textwidth}{@{}p{0.34\textwidth}X@{}}
Arquitecto: & \arquitecto\\[0.15cm]
Líder de Desarrollo: & \liderdesarrollo\\[0.15cm]
Responsable de Calidad: & \responsablecalidad\\[0.15cm]
Responsable de Documentación: & \responsabledocumentacion\\
\end{tabularx}
\vfill
\begin{center}
\textbf{Período académico:} \periodoacademico\par
\medskip
\textbf{Fecha de entrega:} \fechadeentrega\par
\bigskip
\textbf{Repositorio del proyecto:}\par
\smallskip
\mbox{\url{\repositoriopfc}}
\end{center}
\end{titlepage}
}

\tableofcontents
\clearpage

% =====================================================
% Apartados 2-17: editar en docs/secciones/
% La sección 1 corresponde a la portada.
% =====================================================
\setcounter{section}{1}
\seccion{02-resumen}
\seccion{03-contexto}
\seccion{04-objetivos}
\seccion{05-relacionados}
\seccion{06-requisitos}
\seccion{07-justificacion}
\seccion{08-arquitectura}
\seccion{09-comunicacion}
\seccion{10-cap-pacelc}
\seccion{11-adr}
\seccion{12-riesgos}
\seccion{13-plan}
\seccion{14-esqueleto}
\seccion{15-aportes}
\seccion{16-ia}
\seccion{17-referencias}

\clearpage
\appendix
\seccion{18-evidencias}
\clearpage
\seccion{19-transcripciones}

\end{document}
